% SST-CANON-v0.8.1-Appendix.tex
% Supplemental appendix module for Swirl-String-Theory Canon-v0.8.1
% Intended to be included from the main canon after \appendix.
%


\section{Explicit Topological Mass Closures}
\label{sec:appendix_explicit_mass}

\subsection{Purpose and status}
This appendix supplies an explicit closure layer for the mass program without replacing the canonical master equation of Eq.~\eqref{eq:sst_master_equation}. Its role is to make concrete those invariant choices that were historically used in pre-v0.8 layers whenever actual particle-mass calculations were attempted.

\textbf{[DERIVED]} Canon rule: the abstract kernel \(\Xi_K\) of the main text may be specialized to more explicit invariant combinations provided the resulting branch remains dimensionless, benchmarkable, and subordinate to the topology-first hierarchy.

\textbf{[CRITICAL NOTE]} Nothing in this appendix upgrades the explicit kernel to a theorem of first-principles filament dynamics. The present role is closure discipline and reproducible benchmarking.

\subsection{Explicit invariant branch of the canonical kernel}
A conservative explicit specialization is obtained by resolving the topological kernel into ropelength, braid complexity, and Seifert-genus data:
\begin{align}
    \Xi_K^{(\mathrm{exp})}
    =
    \left[
        \alpha_{\mathcal{L}}
        \frac{\mathcal{L}_{\mathrm{tot}}(K)}{\mathcal{L}_{\mathrm{tot}}(3_1)}
        +
        \alpha_b \bigl(b(K)-1\bigr)
        +
        \alpha_g\, g(K)
    \right]
    \varphi^{-2k(K)}.
    \label{eq:explicit_kernel}
\end{align}
Here:
\begin{itemize}
    \item \(\mathcal{L}_{\mathrm{tot}}(K)\) is the total ropelength-like invariant used for the chosen embedding or ideal representative,
    \item \(b(K)\) is the braid index,
    \item \(g(K)\) is the Seifert genus,
    \item \(k(K)\) is the canonical layer / dressing index already used in the main text.
\end{itemize}

\textbf{[ORTHODOX]} The invariants \(\mathcal{L}_{\mathrm{tot}}(K)\), \(b(K)\), and \(g(K)\) are standard topological or geometric knot descriptors.

\textbf{[DERIVED]} Equation~\eqref{eq:explicit_kernel} is dimensionless and therefore admissible as a branch of the canonical kernel.

Substituting Eq.~\eqref{eq:explicit_kernel} into Eq.~\eqref{eq:sst_master_equation} yields the explicit benchmark mass branch
\begin{align}
    M_K^{(\mathrm{exp})}(x)
    =
    \rhoM(x)\,\rc^3\,\Pi_K\,\Xi_K^{(\mathrm{exp})}\,\SwirlClock(x)^{-2}.
    \label{eq:explicit_mass_branch}
\end{align}

\subsection{Baryon benchmark closure}
For composite baryons built from constituent topology classes \(K_u\) and \(K_d\), a conservative additive benchmark branch is
\begin{align}
    M_B^{(\mathrm{add})}
    =
    \Pi_B\,M_0
    \left(
        n_u\,\Xi_{K_u}^{(\mathrm{exp})}
        +
        n_d\,\Xi_{K_d}^{(\mathrm{exp})}
    \right),
    \label{eq:baryon_additive_branch}
\end{align}
where \(n_u\) and \(n_d\) are constituent counts and \(M_0\) is the common medium scale extracted from the main mass functional.

A retained legacy golden-closure benchmark is
\begin{align}
    M_p^{(\varphi)}
    &=
    \varphi^{-2} 3^{-1/\varphi}\,\bigl(2M_u + M_d\bigr),
    \label{eq:proton_phi_benchmark}\\
    M_n^{(\varphi)}
    &=
    \varphi^{-2} 3^{-1/\varphi}\,\bigl(M_u + 2M_d\bigr).
    \label{eq:neutron_phi_benchmark}
\end{align}

\textbf{[SPECULATIVE]} Equations~\eqref{eq:proton_phi_benchmark}--\eqref{eq:neutron_phi_benchmark} are retained as baryon-sector closure benchmarks only. They are not promoted to universal mass laws.

\subsection{Canonical use policy}
The explicit closures above are intended for:
\begin{enumerate}
    \item theorem-ladder benchmarking,
    \item archived CSV reproduction,
    \item side-by-side comparison of topology branches,
    \item explicit baryon fits.
\end{enumerate}
They are \emph{not} intended to replace the abstract organizing form of the main text.

\section{Effective Quantum Bridge}
\label{sec:appendix_quantum_bridge}

\subsection{Purpose and status}
The main canon states that quantum mechanics should appear as an effective low-energy description of the medium. The present appendix supplies the minimal bridge formulas needed to make that statement mathematically legible without claiming a completed derivation of the full quantum rule set.

\textbf{[ORTHODOX]} The Madelung decomposition and Schr\"odinger-form continuum limit are standard.

\textbf{[DERIVED]} SST may use these formulas as a coarse-grained field representation once a local amplitude density and phase have been specified.

\textbf{[SPECULATIVE]} The identification of the resulting fields with branch-resolved swirl carriers remains an SST interpretation.

\subsection{Amplitude--phase decomposition}
Introduce an effective complex mode field
\begin{align}
    \psi(x,t) = \sqrt{\rho_\psi(x,t)}\,e^{i\theta(x,t)},
    \label{eq:madelung_ansatz_appendix}
\end{align}
with \(\rho_\psi \ge 0\) and real phase \(\theta\). The associated effective transport velocity is
\begin{align}
    \mathbf{u}_\psi
    =
    \frac{\hbar}{m_{\mathrm{eff}}}\nabla\theta.
    \label{eq:effective_phase_velocity}
\end{align}

Assuming a Schr\"odinger-form bridge equation
\begin{align}
    i\hbar \partial_t \psi
    =
    \left(
        -\frac{\hbar^2}{2m_{\mathrm{eff}}}\nabla^2 + V_{\mathrm{eff}}
    \right)\psi,
    \label{eq:schrodinger_bridge_appendix}
\end{align}
one obtains the equivalent real system
\begin{align}
    \partial_t \rho_\psi + \nabla\cdot\bigl(\rho_\psi \mathbf{u}_\psi\bigr) &= 0,
    \label{eq:continuity_bridge_appendix}\\
    \partial_t \theta
    +
    \frac{m_{\mathrm{eff}}}{2}\lVert \mathbf{u}_\psi\rVert^2
    + V_{\mathrm{eff}}
    + Q_{\mathrm{eff}} &= 0,
    \label{eq:hamjac_bridge_appendix}
\end{align}
with effective quantum potential
\begin{align}
    Q_{\mathrm{eff}}
    =
    -\frac{\hbar^2}{2m_{\mathrm{eff}}}
    \frac{\nabla^2 \sqrt{\rho_\psi}}{\sqrt{\rho_\psi}}.
    \label{eq:quantum_potential_appendix}
\end{align}

\textbf{[ORTHODOX]} Equations~\eqref{eq:continuity_bridge_appendix}--\eqref{eq:quantum_potential_appendix} are the standard Madelung system associated with Eq.~\eqref{eq:schrodinger_bridge_appendix} \cite{Schrodinger1926,Madelung1927}.

\subsection{SST reading of the bridge}
The conservative SST interpretation is:
\begin{align}
    \rho_\psi
    &\leftrightarrow
    \text{coarse-grained mode-support density},\\
    \theta
    &\leftrightarrow
    \text{effective circulation phase},\\
    \mathbf{u}_\psi
    &\leftrightarrow
    \text{phase-transport velocity of the low-energy envelope}.
\end{align}

At this level, the branch-resolved atomic sector of Section~\ref{sec:atomic} is read not as a replacement of the orthodox wavefunction, but as a candidate microscopic underpinning for the effective fields \(\rho_\psi\) and \(\theta\).

\subsection{Two-component branch bookkeeping}
A minimal two-branch extension consistent with the \(\sigma=\pm1\) structure of the atomic sector is
\begin{align}
    \Psi
    =
    \begin{pmatrix}
        \psi_{+}\\
        \psi_{-}
    \end{pmatrix},
    \qquad
    \psi_{\pm}
    =
    \sqrt{\rho_{\pm}}\,e^{i\theta_{\pm}},
    \label{eq:two_component_bridge}
\end{align}
with total coarse-grained density
\begin{align}
    \rho_{\mathrm{tot}} = \rho_{+}+\rho_{-}.
\end{align}

\textbf{[SPECULATIVE]} This two-component structure is the minimal SST bridge to spinor-like bookkeeping. It is not yet a first-principles derivation of fermionic spin statistics.

\subsection{Topological reading of the bridge}
The most conservative topology statement retained at appendix level is that nontrivial phase maps may encode nontrivial carrier organization. In particular, the effective phase field may carry nontrivial winding or linking data even when the coarse-grained description is written in wave language. This is the appendix-level reason why the canon treats topology and wave structure as compatible rather than as competing descriptions.

\section{Dark-Sector and Galactic Continuum Law}
\label{sec:appendix_dark_galactic}

\subsection{Purpose and status}
The main canon already retains the dark-knot taxonomy and the radial swirl pressure law. This appendix collects the next phenomenological layer: the galaxy-scale force law and the simplest rotation-curve ansatz inherited from earlier SST development.

\textbf{[ORTHODOX]} The underlying force law is ordinary Euler balance in steady swirl.

\textbf{[DERIVED]} SST reinterprets that balance as a candidate large-scale acceleration source.

\textbf{[SPECULATIVE]} The explicit galaxy-fit profile below is a phenomenological closure ansatz, not a theorem.

\subsection{Effective dark acceleration law}
Starting from the steady radial pressure law of Eq.~\eqref{eq:swirl_pressure_law},
\begin{align}
    \frac{1}{\rhoF}\frac{dp_{\mathrm{swirl}}}{dr}
    =
    \frac{v_\theta(r)^2}{r},
\end{align}
define the effective SST acceleration by
\begin{align}
    a_{\mathrm{SST}}(r)
    :=
    \frac{1}{\rhoF}\frac{dp_{\mathrm{swirl}}}{dr}.
    \label{eq:appendix_dark_accel}
\end{align}
Then
\begin{align}
    a_{\mathrm{SST}}(r)=\frac{v_\theta(r)^2}{r}.
    \label{eq:appendix_dark_accel_v2r}
\end{align}

\textbf{[DERIVED]} Equations~\eqref{eq:appendix_dark_accel}--\eqref{eq:appendix_dark_accel_v2r} are immediate consequences of the retained pressure law.

\subsection{Two-component galactic rotation ansatz}
A retained phenomenological velocity profile is
\begin{align}
    v_{\mathrm{swirl}}(r)
    =
    \frac{C_{\mathrm{core}}}{\sqrt{1+\bigl(r_c^{\mathrm{gal}}/r\bigr)^2}}
    +
    C_{\mathrm{tail}}
    \left(1-e^{-r/r_c^{\mathrm{gal}}}\right),
    \label{eq:galactic_velocity_ansatz}
\end{align}
with corresponding total circular speed
\begin{align}
    v_{\mathrm{tot}}^2(r)
    =
    v_{\mathrm{bar}}^2(r)
    +
    v_{\mathrm{swirl}}^2(r).
    \label{eq:total_circular_speed_appendix}
\end{align}

\textbf{[SPECULATIVE]} Equation~\eqref{eq:galactic_velocity_ansatz} is a phenomenological fit ansatz and should be used only at the appendix / research-track level.

\subsection{Relation to dark-knot taxonomy}
The taxonomy of Eq.~\eqref{eq:dark_knot_taxonomy} classifies low-helicity, low-instability knot sectors as dark or quasi-dark candidates. The present galactic law does \emph{not} assume that all galactic anomalies are produced by localized dark knots alone. Rather, it provides the complementary continuum statement: long-range swirl pressure gradients may produce effective extra acceleration even when the underlying microscopic carriers are topologically sparse or weakly luminous.

\subsection{Use policy}
The galactic law belongs in the canon only as:
\begin{enumerate}
    \item an appendix-level phenomenological sector,
    \item a separate astrophysics / dark-sector track,
    \item a fit target for future rotation-curve studies.
\end{enumerate}
It is deliberately not promoted in v0.8.1 to the same status as the primitive constant chain, the Swirl-Clock, or the delay-selection sector.